% !TeX program = xelatex
%% =====================================================================
%%  พรีแอมเบิลร่วม สำหรับสไลด์ประกอบการสอน บทที่ 1 : ตรรกศาสตร์ (Logic)
%%  รายวิชา 4091606 คณิตศาสตร์สำหรับคอมพิวเตอร์
%%  เรียงพิมพ์ด้วย XeLaTeX + ฟอนต์ TH Sarabun New
%%  ใช้ร่วมกันด้วย % !TeX program = xelatex
%% =====================================================================
%%  พรีแอมเบิลร่วม สำหรับสไลด์ประกอบการสอน บทที่ 1 : ตรรกศาสตร์ (Logic)
%%  รายวิชา 4091606 คณิตศาสตร์สำหรับคอมพิวเตอร์
%%  เรียงพิมพ์ด้วย XeLaTeX + ฟอนต์ TH Sarabun New
%%  ใช้ร่วมกันด้วย % !TeX program = xelatex
%% =====================================================================
%%  พรีแอมเบิลร่วม สำหรับสไลด์ประกอบการสอน บทที่ 1 : ตรรกศาสตร์ (Logic)
%%  รายวิชา 4091606 คณิตศาสตร์สำหรับคอมพิวเตอร์
%%  เรียงพิมพ์ด้วย XeLaTeX + ฟอนต์ TH Sarabun New
%%  ใช้ร่วมกันด้วย % !TeX program = xelatex
%% =====================================================================
%%  พรีแอมเบิลร่วม สำหรับสไลด์ประกอบการสอน บทที่ 1 : ตรรกศาสตร์ (Logic)
%%  รายวิชา 4091606 คณิตศาสตร์สำหรับคอมพิวเตอร์
%%  เรียงพิมพ์ด้วย XeLaTeX + ฟอนต์ TH Sarabun New
%%  ใช้ร่วมกันด้วย \input{_preamble} ที่บรรทัดแรกของไฟล์สไลด์แต่ละสัปดาห์
%% =====================================================================
%% โหมด handout: ไฟล์ไดรเวอร์ที่ \def\HANDOUTMODE{} ก่อน \input{_preamble}
%% จะรวมเลเยอร์ overlay (\pause) ให้แสดงเฉลยครบในหน้าเดียว (ไม่มีหน้าซ้ำ)
\ifdefined\HANDOUTMODE\PassOptionsToClass{handout}{beamer}\fi
\documentclass[aspectratio=169,11pt,xcolor=dvipsnames]{beamer}

\usepackage{fontspec}
\usepackage{amsmath,amssymb,amsthm}
\usepackage{booktabs}
\usepackage{array}
\usepackage{colortbl}
\usepackage{tikz}
\usetikzlibrary{arrows.meta,calc,positioning,shapes.gates.logic.US}

%% ---------- ภาษาไทย ----------
%% beamer ใช้แบบอักษร sans-serif เป็นค่าตั้งต้น จึงต้องตั้งทั้ง main และ sans
%% ให้เป็น TH Sarabun New และใช้ professionalfonts กันไม่ให้ beamer แทนที่ฟอนต์
\usefonttheme{professionalfonts}
\XeTeXlinebreaklocale "th_TH"
\XeTeXlinebreakskip = 0pt plus 1pt
\setmainfont[Script=Thai,Scale=1.25]{TH Sarabun New}
\setsansfont[Script=Thai,Scale=1.25]{TH Sarabun New}
\newfontfamily\thaifont[Script=Thai,Scale=1.25]{TH Sarabun New}

%% ---------- จานสี (ยึดตามตำรา) ----------
\definecolor{navy}{HTML}{1F3A93}
\definecolor{Tgreen}{HTML}{2E8B57}
\definecolor{Fred}{HTML}{C0392B}
\definecolor{gold}{HTML}{B7950B}
\definecolor{lightnavy}{HTML}{EAEEF7}
\definecolor{mint}{HTML}{E8F5E9}
\definecolor{lyellow}{HTML}{FFF6D5}
\definecolor{cream}{HTML}{FFF8E1}
\definecolor{lpink}{HTML}{FDEDEC}

%% ---------- ธีม / รูปลักษณ์ ----------
\usetheme{default}
\setbeamertemplate{navigation symbols}{}
\setbeamercolor{frametitle}{fg=white,bg=navy}
\setbeamercolor{title}{fg=navy}
\setbeamercolor{structure}{fg=navy}
\setbeamercolor{itemize item}{fg=navy}
\setbeamercolor{itemize subitem}{fg=gold}
\setbeamercolor{enumerate item}{fg=navy}
\setbeamercolor{section in toc}{fg=navy}
\setbeamercolor{block title}{bg=navy,fg=white}
\setbeamercolor{block body}{bg=lightnavy,fg=black}
\setbeamercolor{block title example}{bg=Tgreen,fg=white}
\setbeamercolor{block body example}{bg=mint,fg=black}
\setbeamercolor{block title alerted}{bg=Fred,fg=white}
\setbeamercolor{block body alerted}{bg=lpink,fg=black}
\setbeamerfont{frametitle}{series=\bfseries}
\setbeamerfont{title}{series=\bfseries}
\setbeamertemplate{itemize item}{\textbullet}
\setbeamertemplate{itemize subitem}{\textendash}

%% เส้นใต้หัวสไลด์
\setbeamertemplate{frametitle}{%
  \nointerlineskip
  \begin{beamercolorbox}[wd=\paperwidth,ht=2.6ex,dp=1.2ex,leftskip=0.3cm,rightskip=0.3cm]{frametitle}%
    \usebeamerfont{frametitle}\insertframetitle%
  \end{beamercolorbox}%
}

%% footline (กำหนดชื่อหัวเรื่องท้ายสไลด์ผ่าน \footchaptertitle)
\newcommand{\footchaptertitle}{บทที่ 1 ตรรกศาสตร์}
\setbeamertemplate{footline}{%
  \leavevmode%
  \hbox{%
    \begin{beamercolorbox}[wd=.5\paperwidth,ht=2.4ex,dp=1ex,leftskip=0.3cm]{author in head/foot}%
      \color{navy}\footnotesize คณิตศาสตร์สำหรับคอมพิวเตอร์ \,|\, \footchaptertitle%
    \end{beamercolorbox}%
    \begin{beamercolorbox}[wd=.5\paperwidth,ht=2.4ex,dp=1ex,rightskip=0.3cm]{author in head/foot}%
      \hfill\color{navy}\footnotesize \insertframenumber\,/\,\inserttotalframenumber%
    \end{beamercolorbox}%
  }%
}

%% ---------- มาโครช่วย ----------
\newcommand{\TT}{\textcolor{Tgreen}{\textbf{T}}}
\newcommand{\FF}{\textcolor{Fred}{\textbf{F}}}
\newcommand{\eng}[1]{\textcolor{gray}{\small (#1)}}
\renewcommand{\arraystretch}{1.25}

%% หน้าแบ่งหัวข้อ (section)
\AtBeginSection[]{%
  \begin{frame}[plain]
    \vfill
    \centering
    \begin{beamercolorbox}[sep=10pt,center,rounded=true,shadow=true]{title}
      \usebeamerfont{title}\Large\insertsectionhead
    \end{beamercolorbox}
    \vfill
  \end{frame}%
}
 ที่บรรทัดแรกของไฟล์สไลด์แต่ละสัปดาห์
%% =====================================================================
%% โหมด handout: ไฟล์ไดรเวอร์ที่ \def\HANDOUTMODE{} ก่อน % !TeX program = xelatex
%% =====================================================================
%%  พรีแอมเบิลร่วม สำหรับสไลด์ประกอบการสอน บทที่ 1 : ตรรกศาสตร์ (Logic)
%%  รายวิชา 4091606 คณิตศาสตร์สำหรับคอมพิวเตอร์
%%  เรียงพิมพ์ด้วย XeLaTeX + ฟอนต์ TH Sarabun New
%%  ใช้ร่วมกันด้วย \input{_preamble} ที่บรรทัดแรกของไฟล์สไลด์แต่ละสัปดาห์
%% =====================================================================
%% โหมด handout: ไฟล์ไดรเวอร์ที่ \def\HANDOUTMODE{} ก่อน \input{_preamble}
%% จะรวมเลเยอร์ overlay (\pause) ให้แสดงเฉลยครบในหน้าเดียว (ไม่มีหน้าซ้ำ)
\ifdefined\HANDOUTMODE\PassOptionsToClass{handout}{beamer}\fi
\documentclass[aspectratio=169,11pt,xcolor=dvipsnames]{beamer}

\usepackage{fontspec}
\usepackage{amsmath,amssymb,amsthm}
\usepackage{booktabs}
\usepackage{array}
\usepackage{colortbl}
\usepackage{tikz}
\usetikzlibrary{arrows.meta,calc,positioning,shapes.gates.logic.US}

%% ---------- ภาษาไทย ----------
%% beamer ใช้แบบอักษร sans-serif เป็นค่าตั้งต้น จึงต้องตั้งทั้ง main และ sans
%% ให้เป็น TH Sarabun New และใช้ professionalfonts กันไม่ให้ beamer แทนที่ฟอนต์
\usefonttheme{professionalfonts}
\XeTeXlinebreaklocale "th_TH"
\XeTeXlinebreakskip = 0pt plus 1pt
\setmainfont[Script=Thai,Scale=1.25]{TH Sarabun New}
\setsansfont[Script=Thai,Scale=1.25]{TH Sarabun New}
\newfontfamily\thaifont[Script=Thai,Scale=1.25]{TH Sarabun New}

%% ---------- จานสี (ยึดตามตำรา) ----------
\definecolor{navy}{HTML}{1F3A93}
\definecolor{Tgreen}{HTML}{2E8B57}
\definecolor{Fred}{HTML}{C0392B}
\definecolor{gold}{HTML}{B7950B}
\definecolor{lightnavy}{HTML}{EAEEF7}
\definecolor{mint}{HTML}{E8F5E9}
\definecolor{lyellow}{HTML}{FFF6D5}
\definecolor{cream}{HTML}{FFF8E1}
\definecolor{lpink}{HTML}{FDEDEC}

%% ---------- ธีม / รูปลักษณ์ ----------
\usetheme{default}
\setbeamertemplate{navigation symbols}{}
\setbeamercolor{frametitle}{fg=white,bg=navy}
\setbeamercolor{title}{fg=navy}
\setbeamercolor{structure}{fg=navy}
\setbeamercolor{itemize item}{fg=navy}
\setbeamercolor{itemize subitem}{fg=gold}
\setbeamercolor{enumerate item}{fg=navy}
\setbeamercolor{section in toc}{fg=navy}
\setbeamercolor{block title}{bg=navy,fg=white}
\setbeamercolor{block body}{bg=lightnavy,fg=black}
\setbeamercolor{block title example}{bg=Tgreen,fg=white}
\setbeamercolor{block body example}{bg=mint,fg=black}
\setbeamercolor{block title alerted}{bg=Fred,fg=white}
\setbeamercolor{block body alerted}{bg=lpink,fg=black}
\setbeamerfont{frametitle}{series=\bfseries}
\setbeamerfont{title}{series=\bfseries}
\setbeamertemplate{itemize item}{\textbullet}
\setbeamertemplate{itemize subitem}{\textendash}

%% เส้นใต้หัวสไลด์
\setbeamertemplate{frametitle}{%
  \nointerlineskip
  \begin{beamercolorbox}[wd=\paperwidth,ht=2.6ex,dp=1.2ex,leftskip=0.3cm,rightskip=0.3cm]{frametitle}%
    \usebeamerfont{frametitle}\insertframetitle%
  \end{beamercolorbox}%
}

%% footline (กำหนดชื่อหัวเรื่องท้ายสไลด์ผ่าน \footchaptertitle)
\newcommand{\footchaptertitle}{บทที่ 1 ตรรกศาสตร์}
\setbeamertemplate{footline}{%
  \leavevmode%
  \hbox{%
    \begin{beamercolorbox}[wd=.5\paperwidth,ht=2.4ex,dp=1ex,leftskip=0.3cm]{author in head/foot}%
      \color{navy}\footnotesize คณิตศาสตร์สำหรับคอมพิวเตอร์ \,|\, \footchaptertitle%
    \end{beamercolorbox}%
    \begin{beamercolorbox}[wd=.5\paperwidth,ht=2.4ex,dp=1ex,rightskip=0.3cm]{author in head/foot}%
      \hfill\color{navy}\footnotesize \insertframenumber\,/\,\inserttotalframenumber%
    \end{beamercolorbox}%
  }%
}

%% ---------- มาโครช่วย ----------
\newcommand{\TT}{\textcolor{Tgreen}{\textbf{T}}}
\newcommand{\FF}{\textcolor{Fred}{\textbf{F}}}
\newcommand{\eng}[1]{\textcolor{gray}{\small (#1)}}
\renewcommand{\arraystretch}{1.25}

%% หน้าแบ่งหัวข้อ (section)
\AtBeginSection[]{%
  \begin{frame}[plain]
    \vfill
    \centering
    \begin{beamercolorbox}[sep=10pt,center,rounded=true,shadow=true]{title}
      \usebeamerfont{title}\Large\insertsectionhead
    \end{beamercolorbox}
    \vfill
  \end{frame}%
}

%% จะรวมเลเยอร์ overlay (\pause) ให้แสดงเฉลยครบในหน้าเดียว (ไม่มีหน้าซ้ำ)
\ifdefined\HANDOUTMODE\PassOptionsToClass{handout}{beamer}\fi
\documentclass[aspectratio=169,11pt,xcolor=dvipsnames]{beamer}

\usepackage{fontspec}
\usepackage{amsmath,amssymb,amsthm}
\usepackage{booktabs}
\usepackage{array}
\usepackage{colortbl}
\usepackage{tikz}
\usetikzlibrary{arrows.meta,calc,positioning,shapes.gates.logic.US}

%% ---------- ภาษาไทย ----------
%% beamer ใช้แบบอักษร sans-serif เป็นค่าตั้งต้น จึงต้องตั้งทั้ง main และ sans
%% ให้เป็น TH Sarabun New และใช้ professionalfonts กันไม่ให้ beamer แทนที่ฟอนต์
\usefonttheme{professionalfonts}
\XeTeXlinebreaklocale "th_TH"
\XeTeXlinebreakskip = 0pt plus 1pt
\setmainfont[Script=Thai,Scale=1.25]{TH Sarabun New}
\setsansfont[Script=Thai,Scale=1.25]{TH Sarabun New}
\newfontfamily\thaifont[Script=Thai,Scale=1.25]{TH Sarabun New}

%% ---------- จานสี (ยึดตามตำรา) ----------
\definecolor{navy}{HTML}{1F3A93}
\definecolor{Tgreen}{HTML}{2E8B57}
\definecolor{Fred}{HTML}{C0392B}
\definecolor{gold}{HTML}{B7950B}
\definecolor{lightnavy}{HTML}{EAEEF7}
\definecolor{mint}{HTML}{E8F5E9}
\definecolor{lyellow}{HTML}{FFF6D5}
\definecolor{cream}{HTML}{FFF8E1}
\definecolor{lpink}{HTML}{FDEDEC}

%% ---------- ธีม / รูปลักษณ์ ----------
\usetheme{default}
\setbeamertemplate{navigation symbols}{}
\setbeamercolor{frametitle}{fg=white,bg=navy}
\setbeamercolor{title}{fg=navy}
\setbeamercolor{structure}{fg=navy}
\setbeamercolor{itemize item}{fg=navy}
\setbeamercolor{itemize subitem}{fg=gold}
\setbeamercolor{enumerate item}{fg=navy}
\setbeamercolor{section in toc}{fg=navy}
\setbeamercolor{block title}{bg=navy,fg=white}
\setbeamercolor{block body}{bg=lightnavy,fg=black}
\setbeamercolor{block title example}{bg=Tgreen,fg=white}
\setbeamercolor{block body example}{bg=mint,fg=black}
\setbeamercolor{block title alerted}{bg=Fred,fg=white}
\setbeamercolor{block body alerted}{bg=lpink,fg=black}
\setbeamerfont{frametitle}{series=\bfseries}
\setbeamerfont{title}{series=\bfseries}
\setbeamertemplate{itemize item}{\textbullet}
\setbeamertemplate{itemize subitem}{\textendash}

%% เส้นใต้หัวสไลด์
\setbeamertemplate{frametitle}{%
  \nointerlineskip
  \begin{beamercolorbox}[wd=\paperwidth,ht=2.6ex,dp=1.2ex,leftskip=0.3cm,rightskip=0.3cm]{frametitle}%
    \usebeamerfont{frametitle}\insertframetitle%
  \end{beamercolorbox}%
}

%% footline (กำหนดชื่อหัวเรื่องท้ายสไลด์ผ่าน \footchaptertitle)
\newcommand{\footchaptertitle}{บทที่ 1 ตรรกศาสตร์}
\setbeamertemplate{footline}{%
  \leavevmode%
  \hbox{%
    \begin{beamercolorbox}[wd=.5\paperwidth,ht=2.4ex,dp=1ex,leftskip=0.3cm]{author in head/foot}%
      \color{navy}\footnotesize คณิตศาสตร์สำหรับคอมพิวเตอร์ \,|\, \footchaptertitle%
    \end{beamercolorbox}%
    \begin{beamercolorbox}[wd=.5\paperwidth,ht=2.4ex,dp=1ex,rightskip=0.3cm]{author in head/foot}%
      \hfill\color{navy}\footnotesize \insertframenumber\,/\,\inserttotalframenumber%
    \end{beamercolorbox}%
  }%
}

%% ---------- มาโครช่วย ----------
\newcommand{\TT}{\textcolor{Tgreen}{\textbf{T}}}
\newcommand{\FF}{\textcolor{Fred}{\textbf{F}}}
\newcommand{\eng}[1]{\textcolor{gray}{\small (#1)}}
\renewcommand{\arraystretch}{1.25}

%% หน้าแบ่งหัวข้อ (section)
\AtBeginSection[]{%
  \begin{frame}[plain]
    \vfill
    \centering
    \begin{beamercolorbox}[sep=10pt,center,rounded=true,shadow=true]{title}
      \usebeamerfont{title}\Large\insertsectionhead
    \end{beamercolorbox}
    \vfill
  \end{frame}%
}
 ที่บรรทัดแรกของไฟล์สไลด์แต่ละสัปดาห์
%% =====================================================================
%% โหมด handout: ไฟล์ไดรเวอร์ที่ \def\HANDOUTMODE{} ก่อน % !TeX program = xelatex
%% =====================================================================
%%  พรีแอมเบิลร่วม สำหรับสไลด์ประกอบการสอน บทที่ 1 : ตรรกศาสตร์ (Logic)
%%  รายวิชา 4091606 คณิตศาสตร์สำหรับคอมพิวเตอร์
%%  เรียงพิมพ์ด้วย XeLaTeX + ฟอนต์ TH Sarabun New
%%  ใช้ร่วมกันด้วย % !TeX program = xelatex
%% =====================================================================
%%  พรีแอมเบิลร่วม สำหรับสไลด์ประกอบการสอน บทที่ 1 : ตรรกศาสตร์ (Logic)
%%  รายวิชา 4091606 คณิตศาสตร์สำหรับคอมพิวเตอร์
%%  เรียงพิมพ์ด้วย XeLaTeX + ฟอนต์ TH Sarabun New
%%  ใช้ร่วมกันด้วย \input{_preamble} ที่บรรทัดแรกของไฟล์สไลด์แต่ละสัปดาห์
%% =====================================================================
%% โหมด handout: ไฟล์ไดรเวอร์ที่ \def\HANDOUTMODE{} ก่อน \input{_preamble}
%% จะรวมเลเยอร์ overlay (\pause) ให้แสดงเฉลยครบในหน้าเดียว (ไม่มีหน้าซ้ำ)
\ifdefined\HANDOUTMODE\PassOptionsToClass{handout}{beamer}\fi
\documentclass[aspectratio=169,11pt,xcolor=dvipsnames]{beamer}

\usepackage{fontspec}
\usepackage{amsmath,amssymb,amsthm}
\usepackage{booktabs}
\usepackage{array}
\usepackage{colortbl}
\usepackage{tikz}
\usetikzlibrary{arrows.meta,calc,positioning,shapes.gates.logic.US}

%% ---------- ภาษาไทย ----------
%% beamer ใช้แบบอักษร sans-serif เป็นค่าตั้งต้น จึงต้องตั้งทั้ง main และ sans
%% ให้เป็น TH Sarabun New และใช้ professionalfonts กันไม่ให้ beamer แทนที่ฟอนต์
\usefonttheme{professionalfonts}
\XeTeXlinebreaklocale "th_TH"
\XeTeXlinebreakskip = 0pt plus 1pt
\setmainfont[Script=Thai,Scale=1.25]{TH Sarabun New}
\setsansfont[Script=Thai,Scale=1.25]{TH Sarabun New}
\newfontfamily\thaifont[Script=Thai,Scale=1.25]{TH Sarabun New}

%% ---------- จานสี (ยึดตามตำรา) ----------
\definecolor{navy}{HTML}{1F3A93}
\definecolor{Tgreen}{HTML}{2E8B57}
\definecolor{Fred}{HTML}{C0392B}
\definecolor{gold}{HTML}{B7950B}
\definecolor{lightnavy}{HTML}{EAEEF7}
\definecolor{mint}{HTML}{E8F5E9}
\definecolor{lyellow}{HTML}{FFF6D5}
\definecolor{cream}{HTML}{FFF8E1}
\definecolor{lpink}{HTML}{FDEDEC}

%% ---------- ธีม / รูปลักษณ์ ----------
\usetheme{default}
\setbeamertemplate{navigation symbols}{}
\setbeamercolor{frametitle}{fg=white,bg=navy}
\setbeamercolor{title}{fg=navy}
\setbeamercolor{structure}{fg=navy}
\setbeamercolor{itemize item}{fg=navy}
\setbeamercolor{itemize subitem}{fg=gold}
\setbeamercolor{enumerate item}{fg=navy}
\setbeamercolor{section in toc}{fg=navy}
\setbeamercolor{block title}{bg=navy,fg=white}
\setbeamercolor{block body}{bg=lightnavy,fg=black}
\setbeamercolor{block title example}{bg=Tgreen,fg=white}
\setbeamercolor{block body example}{bg=mint,fg=black}
\setbeamercolor{block title alerted}{bg=Fred,fg=white}
\setbeamercolor{block body alerted}{bg=lpink,fg=black}
\setbeamerfont{frametitle}{series=\bfseries}
\setbeamerfont{title}{series=\bfseries}
\setbeamertemplate{itemize item}{\textbullet}
\setbeamertemplate{itemize subitem}{\textendash}

%% เส้นใต้หัวสไลด์
\setbeamertemplate{frametitle}{%
  \nointerlineskip
  \begin{beamercolorbox}[wd=\paperwidth,ht=2.6ex,dp=1.2ex,leftskip=0.3cm,rightskip=0.3cm]{frametitle}%
    \usebeamerfont{frametitle}\insertframetitle%
  \end{beamercolorbox}%
}

%% footline (กำหนดชื่อหัวเรื่องท้ายสไลด์ผ่าน \footchaptertitle)
\newcommand{\footchaptertitle}{บทที่ 1 ตรรกศาสตร์}
\setbeamertemplate{footline}{%
  \leavevmode%
  \hbox{%
    \begin{beamercolorbox}[wd=.5\paperwidth,ht=2.4ex,dp=1ex,leftskip=0.3cm]{author in head/foot}%
      \color{navy}\footnotesize คณิตศาสตร์สำหรับคอมพิวเตอร์ \,|\, \footchaptertitle%
    \end{beamercolorbox}%
    \begin{beamercolorbox}[wd=.5\paperwidth,ht=2.4ex,dp=1ex,rightskip=0.3cm]{author in head/foot}%
      \hfill\color{navy}\footnotesize \insertframenumber\,/\,\inserttotalframenumber%
    \end{beamercolorbox}%
  }%
}

%% ---------- มาโครช่วย ----------
\newcommand{\TT}{\textcolor{Tgreen}{\textbf{T}}}
\newcommand{\FF}{\textcolor{Fred}{\textbf{F}}}
\newcommand{\eng}[1]{\textcolor{gray}{\small (#1)}}
\renewcommand{\arraystretch}{1.25}

%% หน้าแบ่งหัวข้อ (section)
\AtBeginSection[]{%
  \begin{frame}[plain]
    \vfill
    \centering
    \begin{beamercolorbox}[sep=10pt,center,rounded=true,shadow=true]{title}
      \usebeamerfont{title}\Large\insertsectionhead
    \end{beamercolorbox}
    \vfill
  \end{frame}%
}
 ที่บรรทัดแรกของไฟล์สไลด์แต่ละสัปดาห์
%% =====================================================================
%% โหมด handout: ไฟล์ไดรเวอร์ที่ \def\HANDOUTMODE{} ก่อน % !TeX program = xelatex
%% =====================================================================
%%  พรีแอมเบิลร่วม สำหรับสไลด์ประกอบการสอน บทที่ 1 : ตรรกศาสตร์ (Logic)
%%  รายวิชา 4091606 คณิตศาสตร์สำหรับคอมพิวเตอร์
%%  เรียงพิมพ์ด้วย XeLaTeX + ฟอนต์ TH Sarabun New
%%  ใช้ร่วมกันด้วย \input{_preamble} ที่บรรทัดแรกของไฟล์สไลด์แต่ละสัปดาห์
%% =====================================================================
%% โหมด handout: ไฟล์ไดรเวอร์ที่ \def\HANDOUTMODE{} ก่อน \input{_preamble}
%% จะรวมเลเยอร์ overlay (\pause) ให้แสดงเฉลยครบในหน้าเดียว (ไม่มีหน้าซ้ำ)
\ifdefined\HANDOUTMODE\PassOptionsToClass{handout}{beamer}\fi
\documentclass[aspectratio=169,11pt,xcolor=dvipsnames]{beamer}

\usepackage{fontspec}
\usepackage{amsmath,amssymb,amsthm}
\usepackage{booktabs}
\usepackage{array}
\usepackage{colortbl}
\usepackage{tikz}
\usetikzlibrary{arrows.meta,calc,positioning,shapes.gates.logic.US}

%% ---------- ภาษาไทย ----------
%% beamer ใช้แบบอักษร sans-serif เป็นค่าตั้งต้น จึงต้องตั้งทั้ง main และ sans
%% ให้เป็น TH Sarabun New และใช้ professionalfonts กันไม่ให้ beamer แทนที่ฟอนต์
\usefonttheme{professionalfonts}
\XeTeXlinebreaklocale "th_TH"
\XeTeXlinebreakskip = 0pt plus 1pt
\setmainfont[Script=Thai,Scale=1.25]{TH Sarabun New}
\setsansfont[Script=Thai,Scale=1.25]{TH Sarabun New}
\newfontfamily\thaifont[Script=Thai,Scale=1.25]{TH Sarabun New}

%% ---------- จานสี (ยึดตามตำรา) ----------
\definecolor{navy}{HTML}{1F3A93}
\definecolor{Tgreen}{HTML}{2E8B57}
\definecolor{Fred}{HTML}{C0392B}
\definecolor{gold}{HTML}{B7950B}
\definecolor{lightnavy}{HTML}{EAEEF7}
\definecolor{mint}{HTML}{E8F5E9}
\definecolor{lyellow}{HTML}{FFF6D5}
\definecolor{cream}{HTML}{FFF8E1}
\definecolor{lpink}{HTML}{FDEDEC}

%% ---------- ธีม / รูปลักษณ์ ----------
\usetheme{default}
\setbeamertemplate{navigation symbols}{}
\setbeamercolor{frametitle}{fg=white,bg=navy}
\setbeamercolor{title}{fg=navy}
\setbeamercolor{structure}{fg=navy}
\setbeamercolor{itemize item}{fg=navy}
\setbeamercolor{itemize subitem}{fg=gold}
\setbeamercolor{enumerate item}{fg=navy}
\setbeamercolor{section in toc}{fg=navy}
\setbeamercolor{block title}{bg=navy,fg=white}
\setbeamercolor{block body}{bg=lightnavy,fg=black}
\setbeamercolor{block title example}{bg=Tgreen,fg=white}
\setbeamercolor{block body example}{bg=mint,fg=black}
\setbeamercolor{block title alerted}{bg=Fred,fg=white}
\setbeamercolor{block body alerted}{bg=lpink,fg=black}
\setbeamerfont{frametitle}{series=\bfseries}
\setbeamerfont{title}{series=\bfseries}
\setbeamertemplate{itemize item}{\textbullet}
\setbeamertemplate{itemize subitem}{\textendash}

%% เส้นใต้หัวสไลด์
\setbeamertemplate{frametitle}{%
  \nointerlineskip
  \begin{beamercolorbox}[wd=\paperwidth,ht=2.6ex,dp=1.2ex,leftskip=0.3cm,rightskip=0.3cm]{frametitle}%
    \usebeamerfont{frametitle}\insertframetitle%
  \end{beamercolorbox}%
}

%% footline (กำหนดชื่อหัวเรื่องท้ายสไลด์ผ่าน \footchaptertitle)
\newcommand{\footchaptertitle}{บทที่ 1 ตรรกศาสตร์}
\setbeamertemplate{footline}{%
  \leavevmode%
  \hbox{%
    \begin{beamercolorbox}[wd=.5\paperwidth,ht=2.4ex,dp=1ex,leftskip=0.3cm]{author in head/foot}%
      \color{navy}\footnotesize คณิตศาสตร์สำหรับคอมพิวเตอร์ \,|\, \footchaptertitle%
    \end{beamercolorbox}%
    \begin{beamercolorbox}[wd=.5\paperwidth,ht=2.4ex,dp=1ex,rightskip=0.3cm]{author in head/foot}%
      \hfill\color{navy}\footnotesize \insertframenumber\,/\,\inserttotalframenumber%
    \end{beamercolorbox}%
  }%
}

%% ---------- มาโครช่วย ----------
\newcommand{\TT}{\textcolor{Tgreen}{\textbf{T}}}
\newcommand{\FF}{\textcolor{Fred}{\textbf{F}}}
\newcommand{\eng}[1]{\textcolor{gray}{\small (#1)}}
\renewcommand{\arraystretch}{1.25}

%% หน้าแบ่งหัวข้อ (section)
\AtBeginSection[]{%
  \begin{frame}[plain]
    \vfill
    \centering
    \begin{beamercolorbox}[sep=10pt,center,rounded=true,shadow=true]{title}
      \usebeamerfont{title}\Large\insertsectionhead
    \end{beamercolorbox}
    \vfill
  \end{frame}%
}

%% จะรวมเลเยอร์ overlay (\pause) ให้แสดงเฉลยครบในหน้าเดียว (ไม่มีหน้าซ้ำ)
\ifdefined\HANDOUTMODE\PassOptionsToClass{handout}{beamer}\fi
\documentclass[aspectratio=169,11pt,xcolor=dvipsnames]{beamer}

\usepackage{fontspec}
\usepackage{amsmath,amssymb,amsthm}
\usepackage{booktabs}
\usepackage{array}
\usepackage{colortbl}
\usepackage{tikz}
\usetikzlibrary{arrows.meta,calc,positioning,shapes.gates.logic.US}

%% ---------- ภาษาไทย ----------
%% beamer ใช้แบบอักษร sans-serif เป็นค่าตั้งต้น จึงต้องตั้งทั้ง main และ sans
%% ให้เป็น TH Sarabun New และใช้ professionalfonts กันไม่ให้ beamer แทนที่ฟอนต์
\usefonttheme{professionalfonts}
\XeTeXlinebreaklocale "th_TH"
\XeTeXlinebreakskip = 0pt plus 1pt
\setmainfont[Script=Thai,Scale=1.25]{TH Sarabun New}
\setsansfont[Script=Thai,Scale=1.25]{TH Sarabun New}
\newfontfamily\thaifont[Script=Thai,Scale=1.25]{TH Sarabun New}

%% ---------- จานสี (ยึดตามตำรา) ----------
\definecolor{navy}{HTML}{1F3A93}
\definecolor{Tgreen}{HTML}{2E8B57}
\definecolor{Fred}{HTML}{C0392B}
\definecolor{gold}{HTML}{B7950B}
\definecolor{lightnavy}{HTML}{EAEEF7}
\definecolor{mint}{HTML}{E8F5E9}
\definecolor{lyellow}{HTML}{FFF6D5}
\definecolor{cream}{HTML}{FFF8E1}
\definecolor{lpink}{HTML}{FDEDEC}

%% ---------- ธีม / รูปลักษณ์ ----------
\usetheme{default}
\setbeamertemplate{navigation symbols}{}
\setbeamercolor{frametitle}{fg=white,bg=navy}
\setbeamercolor{title}{fg=navy}
\setbeamercolor{structure}{fg=navy}
\setbeamercolor{itemize item}{fg=navy}
\setbeamercolor{itemize subitem}{fg=gold}
\setbeamercolor{enumerate item}{fg=navy}
\setbeamercolor{section in toc}{fg=navy}
\setbeamercolor{block title}{bg=navy,fg=white}
\setbeamercolor{block body}{bg=lightnavy,fg=black}
\setbeamercolor{block title example}{bg=Tgreen,fg=white}
\setbeamercolor{block body example}{bg=mint,fg=black}
\setbeamercolor{block title alerted}{bg=Fred,fg=white}
\setbeamercolor{block body alerted}{bg=lpink,fg=black}
\setbeamerfont{frametitle}{series=\bfseries}
\setbeamerfont{title}{series=\bfseries}
\setbeamertemplate{itemize item}{\textbullet}
\setbeamertemplate{itemize subitem}{\textendash}

%% เส้นใต้หัวสไลด์
\setbeamertemplate{frametitle}{%
  \nointerlineskip
  \begin{beamercolorbox}[wd=\paperwidth,ht=2.6ex,dp=1.2ex,leftskip=0.3cm,rightskip=0.3cm]{frametitle}%
    \usebeamerfont{frametitle}\insertframetitle%
  \end{beamercolorbox}%
}

%% footline (กำหนดชื่อหัวเรื่องท้ายสไลด์ผ่าน \footchaptertitle)
\newcommand{\footchaptertitle}{บทที่ 1 ตรรกศาสตร์}
\setbeamertemplate{footline}{%
  \leavevmode%
  \hbox{%
    \begin{beamercolorbox}[wd=.5\paperwidth,ht=2.4ex,dp=1ex,leftskip=0.3cm]{author in head/foot}%
      \color{navy}\footnotesize คณิตศาสตร์สำหรับคอมพิวเตอร์ \,|\, \footchaptertitle%
    \end{beamercolorbox}%
    \begin{beamercolorbox}[wd=.5\paperwidth,ht=2.4ex,dp=1ex,rightskip=0.3cm]{author in head/foot}%
      \hfill\color{navy}\footnotesize \insertframenumber\,/\,\inserttotalframenumber%
    \end{beamercolorbox}%
  }%
}

%% ---------- มาโครช่วย ----------
\newcommand{\TT}{\textcolor{Tgreen}{\textbf{T}}}
\newcommand{\FF}{\textcolor{Fred}{\textbf{F}}}
\newcommand{\eng}[1]{\textcolor{gray}{\small (#1)}}
\renewcommand{\arraystretch}{1.25}

%% หน้าแบ่งหัวข้อ (section)
\AtBeginSection[]{%
  \begin{frame}[plain]
    \vfill
    \centering
    \begin{beamercolorbox}[sep=10pt,center,rounded=true,shadow=true]{title}
      \usebeamerfont{title}\Large\insertsectionhead
    \end{beamercolorbox}
    \vfill
  \end{frame}%
}

%% จะรวมเลเยอร์ overlay (\pause) ให้แสดงเฉลยครบในหน้าเดียว (ไม่มีหน้าซ้ำ)
\ifdefined\HANDOUTMODE\PassOptionsToClass{handout}{beamer}\fi
\documentclass[aspectratio=169,11pt,xcolor=dvipsnames]{beamer}

\usepackage{fontspec}
\usepackage{amsmath,amssymb,amsthm}
\usepackage{booktabs}
\usepackage{array}
\usepackage{colortbl}
\usepackage{tikz}
\usetikzlibrary{arrows.meta,calc,positioning,shapes.gates.logic.US}

%% ---------- ภาษาไทย ----------
%% beamer ใช้แบบอักษร sans-serif เป็นค่าตั้งต้น จึงต้องตั้งทั้ง main และ sans
%% ให้เป็น TH Sarabun New และใช้ professionalfonts กันไม่ให้ beamer แทนที่ฟอนต์
\usefonttheme{professionalfonts}
\XeTeXlinebreaklocale "th_TH"
\XeTeXlinebreakskip = 0pt plus 1pt
\setmainfont[Script=Thai,Scale=1.25]{TH Sarabun New}
\setsansfont[Script=Thai,Scale=1.25]{TH Sarabun New}
\newfontfamily\thaifont[Script=Thai,Scale=1.25]{TH Sarabun New}

%% ---------- จานสี (ยึดตามตำรา) ----------
\definecolor{navy}{HTML}{1F3A93}
\definecolor{Tgreen}{HTML}{2E8B57}
\definecolor{Fred}{HTML}{C0392B}
\definecolor{gold}{HTML}{B7950B}
\definecolor{lightnavy}{HTML}{EAEEF7}
\definecolor{mint}{HTML}{E8F5E9}
\definecolor{lyellow}{HTML}{FFF6D5}
\definecolor{cream}{HTML}{FFF8E1}
\definecolor{lpink}{HTML}{FDEDEC}

%% ---------- ธีม / รูปลักษณ์ ----------
\usetheme{default}
\setbeamertemplate{navigation symbols}{}
\setbeamercolor{frametitle}{fg=white,bg=navy}
\setbeamercolor{title}{fg=navy}
\setbeamercolor{structure}{fg=navy}
\setbeamercolor{itemize item}{fg=navy}
\setbeamercolor{itemize subitem}{fg=gold}
\setbeamercolor{enumerate item}{fg=navy}
\setbeamercolor{section in toc}{fg=navy}
\setbeamercolor{block title}{bg=navy,fg=white}
\setbeamercolor{block body}{bg=lightnavy,fg=black}
\setbeamercolor{block title example}{bg=Tgreen,fg=white}
\setbeamercolor{block body example}{bg=mint,fg=black}
\setbeamercolor{block title alerted}{bg=Fred,fg=white}
\setbeamercolor{block body alerted}{bg=lpink,fg=black}
\setbeamerfont{frametitle}{series=\bfseries}
\setbeamerfont{title}{series=\bfseries}
\setbeamertemplate{itemize item}{\textbullet}
\setbeamertemplate{itemize subitem}{\textendash}

%% เส้นใต้หัวสไลด์
\setbeamertemplate{frametitle}{%
  \nointerlineskip
  \begin{beamercolorbox}[wd=\paperwidth,ht=2.6ex,dp=1.2ex,leftskip=0.3cm,rightskip=0.3cm]{frametitle}%
    \usebeamerfont{frametitle}\insertframetitle%
  \end{beamercolorbox}%
}

%% footline (กำหนดชื่อหัวเรื่องท้ายสไลด์ผ่าน \footchaptertitle)
\newcommand{\footchaptertitle}{บทที่ 1 ตรรกศาสตร์}
\setbeamertemplate{footline}{%
  \leavevmode%
  \hbox{%
    \begin{beamercolorbox}[wd=.5\paperwidth,ht=2.4ex,dp=1ex,leftskip=0.3cm]{author in head/foot}%
      \color{navy}\footnotesize คณิตศาสตร์สำหรับคอมพิวเตอร์ \,|\, \footchaptertitle%
    \end{beamercolorbox}%
    \begin{beamercolorbox}[wd=.5\paperwidth,ht=2.4ex,dp=1ex,rightskip=0.3cm]{author in head/foot}%
      \hfill\color{navy}\footnotesize \insertframenumber\,/\,\inserttotalframenumber%
    \end{beamercolorbox}%
  }%
}

%% ---------- มาโครช่วย ----------
\newcommand{\TT}{\textcolor{Tgreen}{\textbf{T}}}
\newcommand{\FF}{\textcolor{Fred}{\textbf{F}}}
\newcommand{\eng}[1]{\textcolor{gray}{\small (#1)}}
\renewcommand{\arraystretch}{1.25}

%% หน้าแบ่งหัวข้อ (section)
\AtBeginSection[]{%
  \begin{frame}[plain]
    \vfill
    \centering
    \begin{beamercolorbox}[sep=10pt,center,rounded=true,shadow=true]{title}
      \usebeamerfont{title}\Large\insertsectionhead
    \end{beamercolorbox}
    \vfill
  \end{frame}%
}
 ที่บรรทัดแรกของไฟล์สไลด์แต่ละสัปดาห์
%% =====================================================================
%% โหมด handout: ไฟล์ไดรเวอร์ที่ \def\HANDOUTMODE{} ก่อน % !TeX program = xelatex
%% =====================================================================
%%  พรีแอมเบิลร่วม สำหรับสไลด์ประกอบการสอน บทที่ 1 : ตรรกศาสตร์ (Logic)
%%  รายวิชา 4091606 คณิตศาสตร์สำหรับคอมพิวเตอร์
%%  เรียงพิมพ์ด้วย XeLaTeX + ฟอนต์ TH Sarabun New
%%  ใช้ร่วมกันด้วย % !TeX program = xelatex
%% =====================================================================
%%  พรีแอมเบิลร่วม สำหรับสไลด์ประกอบการสอน บทที่ 1 : ตรรกศาสตร์ (Logic)
%%  รายวิชา 4091606 คณิตศาสตร์สำหรับคอมพิวเตอร์
%%  เรียงพิมพ์ด้วย XeLaTeX + ฟอนต์ TH Sarabun New
%%  ใช้ร่วมกันด้วย % !TeX program = xelatex
%% =====================================================================
%%  พรีแอมเบิลร่วม สำหรับสไลด์ประกอบการสอน บทที่ 1 : ตรรกศาสตร์ (Logic)
%%  รายวิชา 4091606 คณิตศาสตร์สำหรับคอมพิวเตอร์
%%  เรียงพิมพ์ด้วย XeLaTeX + ฟอนต์ TH Sarabun New
%%  ใช้ร่วมกันด้วย \input{_preamble} ที่บรรทัดแรกของไฟล์สไลด์แต่ละสัปดาห์
%% =====================================================================
%% โหมด handout: ไฟล์ไดรเวอร์ที่ \def\HANDOUTMODE{} ก่อน \input{_preamble}
%% จะรวมเลเยอร์ overlay (\pause) ให้แสดงเฉลยครบในหน้าเดียว (ไม่มีหน้าซ้ำ)
\ifdefined\HANDOUTMODE\PassOptionsToClass{handout}{beamer}\fi
\documentclass[aspectratio=169,11pt,xcolor=dvipsnames]{beamer}

\usepackage{fontspec}
\usepackage{amsmath,amssymb,amsthm}
\usepackage{booktabs}
\usepackage{array}
\usepackage{colortbl}
\usepackage{tikz}
\usetikzlibrary{arrows.meta,calc,positioning,shapes.gates.logic.US}

%% ---------- ภาษาไทย ----------
%% beamer ใช้แบบอักษร sans-serif เป็นค่าตั้งต้น จึงต้องตั้งทั้ง main และ sans
%% ให้เป็น TH Sarabun New และใช้ professionalfonts กันไม่ให้ beamer แทนที่ฟอนต์
\usefonttheme{professionalfonts}
\XeTeXlinebreaklocale "th_TH"
\XeTeXlinebreakskip = 0pt plus 1pt
\setmainfont[Script=Thai,Scale=1.25]{TH Sarabun New}
\setsansfont[Script=Thai,Scale=1.25]{TH Sarabun New}
\newfontfamily\thaifont[Script=Thai,Scale=1.25]{TH Sarabun New}

%% ---------- จานสี (ยึดตามตำรา) ----------
\definecolor{navy}{HTML}{1F3A93}
\definecolor{Tgreen}{HTML}{2E8B57}
\definecolor{Fred}{HTML}{C0392B}
\definecolor{gold}{HTML}{B7950B}
\definecolor{lightnavy}{HTML}{EAEEF7}
\definecolor{mint}{HTML}{E8F5E9}
\definecolor{lyellow}{HTML}{FFF6D5}
\definecolor{cream}{HTML}{FFF8E1}
\definecolor{lpink}{HTML}{FDEDEC}

%% ---------- ธีม / รูปลักษณ์ ----------
\usetheme{default}
\setbeamertemplate{navigation symbols}{}
\setbeamercolor{frametitle}{fg=white,bg=navy}
\setbeamercolor{title}{fg=navy}
\setbeamercolor{structure}{fg=navy}
\setbeamercolor{itemize item}{fg=navy}
\setbeamercolor{itemize subitem}{fg=gold}
\setbeamercolor{enumerate item}{fg=navy}
\setbeamercolor{section in toc}{fg=navy}
\setbeamercolor{block title}{bg=navy,fg=white}
\setbeamercolor{block body}{bg=lightnavy,fg=black}
\setbeamercolor{block title example}{bg=Tgreen,fg=white}
\setbeamercolor{block body example}{bg=mint,fg=black}
\setbeamercolor{block title alerted}{bg=Fred,fg=white}
\setbeamercolor{block body alerted}{bg=lpink,fg=black}
\setbeamerfont{frametitle}{series=\bfseries}
\setbeamerfont{title}{series=\bfseries}
\setbeamertemplate{itemize item}{\textbullet}
\setbeamertemplate{itemize subitem}{\textendash}

%% เส้นใต้หัวสไลด์
\setbeamertemplate{frametitle}{%
  \nointerlineskip
  \begin{beamercolorbox}[wd=\paperwidth,ht=2.6ex,dp=1.2ex,leftskip=0.3cm,rightskip=0.3cm]{frametitle}%
    \usebeamerfont{frametitle}\insertframetitle%
  \end{beamercolorbox}%
}

%% footline (กำหนดชื่อหัวเรื่องท้ายสไลด์ผ่าน \footchaptertitle)
\newcommand{\footchaptertitle}{บทที่ 1 ตรรกศาสตร์}
\setbeamertemplate{footline}{%
  \leavevmode%
  \hbox{%
    \begin{beamercolorbox}[wd=.5\paperwidth,ht=2.4ex,dp=1ex,leftskip=0.3cm]{author in head/foot}%
      \color{navy}\footnotesize คณิตศาสตร์สำหรับคอมพิวเตอร์ \,|\, \footchaptertitle%
    \end{beamercolorbox}%
    \begin{beamercolorbox}[wd=.5\paperwidth,ht=2.4ex,dp=1ex,rightskip=0.3cm]{author in head/foot}%
      \hfill\color{navy}\footnotesize \insertframenumber\,/\,\inserttotalframenumber%
    \end{beamercolorbox}%
  }%
}

%% ---------- มาโครช่วย ----------
\newcommand{\TT}{\textcolor{Tgreen}{\textbf{T}}}
\newcommand{\FF}{\textcolor{Fred}{\textbf{F}}}
\newcommand{\eng}[1]{\textcolor{gray}{\small (#1)}}
\renewcommand{\arraystretch}{1.25}

%% หน้าแบ่งหัวข้อ (section)
\AtBeginSection[]{%
  \begin{frame}[plain]
    \vfill
    \centering
    \begin{beamercolorbox}[sep=10pt,center,rounded=true,shadow=true]{title}
      \usebeamerfont{title}\Large\insertsectionhead
    \end{beamercolorbox}
    \vfill
  \end{frame}%
}
 ที่บรรทัดแรกของไฟล์สไลด์แต่ละสัปดาห์
%% =====================================================================
%% โหมด handout: ไฟล์ไดรเวอร์ที่ \def\HANDOUTMODE{} ก่อน % !TeX program = xelatex
%% =====================================================================
%%  พรีแอมเบิลร่วม สำหรับสไลด์ประกอบการสอน บทที่ 1 : ตรรกศาสตร์ (Logic)
%%  รายวิชา 4091606 คณิตศาสตร์สำหรับคอมพิวเตอร์
%%  เรียงพิมพ์ด้วย XeLaTeX + ฟอนต์ TH Sarabun New
%%  ใช้ร่วมกันด้วย \input{_preamble} ที่บรรทัดแรกของไฟล์สไลด์แต่ละสัปดาห์
%% =====================================================================
%% โหมด handout: ไฟล์ไดรเวอร์ที่ \def\HANDOUTMODE{} ก่อน \input{_preamble}
%% จะรวมเลเยอร์ overlay (\pause) ให้แสดงเฉลยครบในหน้าเดียว (ไม่มีหน้าซ้ำ)
\ifdefined\HANDOUTMODE\PassOptionsToClass{handout}{beamer}\fi
\documentclass[aspectratio=169,11pt,xcolor=dvipsnames]{beamer}

\usepackage{fontspec}
\usepackage{amsmath,amssymb,amsthm}
\usepackage{booktabs}
\usepackage{array}
\usepackage{colortbl}
\usepackage{tikz}
\usetikzlibrary{arrows.meta,calc,positioning,shapes.gates.logic.US}

%% ---------- ภาษาไทย ----------
%% beamer ใช้แบบอักษร sans-serif เป็นค่าตั้งต้น จึงต้องตั้งทั้ง main และ sans
%% ให้เป็น TH Sarabun New และใช้ professionalfonts กันไม่ให้ beamer แทนที่ฟอนต์
\usefonttheme{professionalfonts}
\XeTeXlinebreaklocale "th_TH"
\XeTeXlinebreakskip = 0pt plus 1pt
\setmainfont[Script=Thai,Scale=1.25]{TH Sarabun New}
\setsansfont[Script=Thai,Scale=1.25]{TH Sarabun New}
\newfontfamily\thaifont[Script=Thai,Scale=1.25]{TH Sarabun New}

%% ---------- จานสี (ยึดตามตำรา) ----------
\definecolor{navy}{HTML}{1F3A93}
\definecolor{Tgreen}{HTML}{2E8B57}
\definecolor{Fred}{HTML}{C0392B}
\definecolor{gold}{HTML}{B7950B}
\definecolor{lightnavy}{HTML}{EAEEF7}
\definecolor{mint}{HTML}{E8F5E9}
\definecolor{lyellow}{HTML}{FFF6D5}
\definecolor{cream}{HTML}{FFF8E1}
\definecolor{lpink}{HTML}{FDEDEC}

%% ---------- ธีม / รูปลักษณ์ ----------
\usetheme{default}
\setbeamertemplate{navigation symbols}{}
\setbeamercolor{frametitle}{fg=white,bg=navy}
\setbeamercolor{title}{fg=navy}
\setbeamercolor{structure}{fg=navy}
\setbeamercolor{itemize item}{fg=navy}
\setbeamercolor{itemize subitem}{fg=gold}
\setbeamercolor{enumerate item}{fg=navy}
\setbeamercolor{section in toc}{fg=navy}
\setbeamercolor{block title}{bg=navy,fg=white}
\setbeamercolor{block body}{bg=lightnavy,fg=black}
\setbeamercolor{block title example}{bg=Tgreen,fg=white}
\setbeamercolor{block body example}{bg=mint,fg=black}
\setbeamercolor{block title alerted}{bg=Fred,fg=white}
\setbeamercolor{block body alerted}{bg=lpink,fg=black}
\setbeamerfont{frametitle}{series=\bfseries}
\setbeamerfont{title}{series=\bfseries}
\setbeamertemplate{itemize item}{\textbullet}
\setbeamertemplate{itemize subitem}{\textendash}

%% เส้นใต้หัวสไลด์
\setbeamertemplate{frametitle}{%
  \nointerlineskip
  \begin{beamercolorbox}[wd=\paperwidth,ht=2.6ex,dp=1.2ex,leftskip=0.3cm,rightskip=0.3cm]{frametitle}%
    \usebeamerfont{frametitle}\insertframetitle%
  \end{beamercolorbox}%
}

%% footline (กำหนดชื่อหัวเรื่องท้ายสไลด์ผ่าน \footchaptertitle)
\newcommand{\footchaptertitle}{บทที่ 1 ตรรกศาสตร์}
\setbeamertemplate{footline}{%
  \leavevmode%
  \hbox{%
    \begin{beamercolorbox}[wd=.5\paperwidth,ht=2.4ex,dp=1ex,leftskip=0.3cm]{author in head/foot}%
      \color{navy}\footnotesize คณิตศาสตร์สำหรับคอมพิวเตอร์ \,|\, \footchaptertitle%
    \end{beamercolorbox}%
    \begin{beamercolorbox}[wd=.5\paperwidth,ht=2.4ex,dp=1ex,rightskip=0.3cm]{author in head/foot}%
      \hfill\color{navy}\footnotesize \insertframenumber\,/\,\inserttotalframenumber%
    \end{beamercolorbox}%
  }%
}

%% ---------- มาโครช่วย ----------
\newcommand{\TT}{\textcolor{Tgreen}{\textbf{T}}}
\newcommand{\FF}{\textcolor{Fred}{\textbf{F}}}
\newcommand{\eng}[1]{\textcolor{gray}{\small (#1)}}
\renewcommand{\arraystretch}{1.25}

%% หน้าแบ่งหัวข้อ (section)
\AtBeginSection[]{%
  \begin{frame}[plain]
    \vfill
    \centering
    \begin{beamercolorbox}[sep=10pt,center,rounded=true,shadow=true]{title}
      \usebeamerfont{title}\Large\insertsectionhead
    \end{beamercolorbox}
    \vfill
  \end{frame}%
}

%% จะรวมเลเยอร์ overlay (\pause) ให้แสดงเฉลยครบในหน้าเดียว (ไม่มีหน้าซ้ำ)
\ifdefined\HANDOUTMODE\PassOptionsToClass{handout}{beamer}\fi
\documentclass[aspectratio=169,11pt,xcolor=dvipsnames]{beamer}

\usepackage{fontspec}
\usepackage{amsmath,amssymb,amsthm}
\usepackage{booktabs}
\usepackage{array}
\usepackage{colortbl}
\usepackage{tikz}
\usetikzlibrary{arrows.meta,calc,positioning,shapes.gates.logic.US}

%% ---------- ภาษาไทย ----------
%% beamer ใช้แบบอักษร sans-serif เป็นค่าตั้งต้น จึงต้องตั้งทั้ง main และ sans
%% ให้เป็น TH Sarabun New และใช้ professionalfonts กันไม่ให้ beamer แทนที่ฟอนต์
\usefonttheme{professionalfonts}
\XeTeXlinebreaklocale "th_TH"
\XeTeXlinebreakskip = 0pt plus 1pt
\setmainfont[Script=Thai,Scale=1.25]{TH Sarabun New}
\setsansfont[Script=Thai,Scale=1.25]{TH Sarabun New}
\newfontfamily\thaifont[Script=Thai,Scale=1.25]{TH Sarabun New}

%% ---------- จานสี (ยึดตามตำรา) ----------
\definecolor{navy}{HTML}{1F3A93}
\definecolor{Tgreen}{HTML}{2E8B57}
\definecolor{Fred}{HTML}{C0392B}
\definecolor{gold}{HTML}{B7950B}
\definecolor{lightnavy}{HTML}{EAEEF7}
\definecolor{mint}{HTML}{E8F5E9}
\definecolor{lyellow}{HTML}{FFF6D5}
\definecolor{cream}{HTML}{FFF8E1}
\definecolor{lpink}{HTML}{FDEDEC}

%% ---------- ธีม / รูปลักษณ์ ----------
\usetheme{default}
\setbeamertemplate{navigation symbols}{}
\setbeamercolor{frametitle}{fg=white,bg=navy}
\setbeamercolor{title}{fg=navy}
\setbeamercolor{structure}{fg=navy}
\setbeamercolor{itemize item}{fg=navy}
\setbeamercolor{itemize subitem}{fg=gold}
\setbeamercolor{enumerate item}{fg=navy}
\setbeamercolor{section in toc}{fg=navy}
\setbeamercolor{block title}{bg=navy,fg=white}
\setbeamercolor{block body}{bg=lightnavy,fg=black}
\setbeamercolor{block title example}{bg=Tgreen,fg=white}
\setbeamercolor{block body example}{bg=mint,fg=black}
\setbeamercolor{block title alerted}{bg=Fred,fg=white}
\setbeamercolor{block body alerted}{bg=lpink,fg=black}
\setbeamerfont{frametitle}{series=\bfseries}
\setbeamerfont{title}{series=\bfseries}
\setbeamertemplate{itemize item}{\textbullet}
\setbeamertemplate{itemize subitem}{\textendash}

%% เส้นใต้หัวสไลด์
\setbeamertemplate{frametitle}{%
  \nointerlineskip
  \begin{beamercolorbox}[wd=\paperwidth,ht=2.6ex,dp=1.2ex,leftskip=0.3cm,rightskip=0.3cm]{frametitle}%
    \usebeamerfont{frametitle}\insertframetitle%
  \end{beamercolorbox}%
}

%% footline (กำหนดชื่อหัวเรื่องท้ายสไลด์ผ่าน \footchaptertitle)
\newcommand{\footchaptertitle}{บทที่ 1 ตรรกศาสตร์}
\setbeamertemplate{footline}{%
  \leavevmode%
  \hbox{%
    \begin{beamercolorbox}[wd=.5\paperwidth,ht=2.4ex,dp=1ex,leftskip=0.3cm]{author in head/foot}%
      \color{navy}\footnotesize คณิตศาสตร์สำหรับคอมพิวเตอร์ \,|\, \footchaptertitle%
    \end{beamercolorbox}%
    \begin{beamercolorbox}[wd=.5\paperwidth,ht=2.4ex,dp=1ex,rightskip=0.3cm]{author in head/foot}%
      \hfill\color{navy}\footnotesize \insertframenumber\,/\,\inserttotalframenumber%
    \end{beamercolorbox}%
  }%
}

%% ---------- มาโครช่วย ----------
\newcommand{\TT}{\textcolor{Tgreen}{\textbf{T}}}
\newcommand{\FF}{\textcolor{Fred}{\textbf{F}}}
\newcommand{\eng}[1]{\textcolor{gray}{\small (#1)}}
\renewcommand{\arraystretch}{1.25}

%% หน้าแบ่งหัวข้อ (section)
\AtBeginSection[]{%
  \begin{frame}[plain]
    \vfill
    \centering
    \begin{beamercolorbox}[sep=10pt,center,rounded=true,shadow=true]{title}
      \usebeamerfont{title}\Large\insertsectionhead
    \end{beamercolorbox}
    \vfill
  \end{frame}%
}
 ที่บรรทัดแรกของไฟล์สไลด์แต่ละสัปดาห์
%% =====================================================================
%% โหมด handout: ไฟล์ไดรเวอร์ที่ \def\HANDOUTMODE{} ก่อน % !TeX program = xelatex
%% =====================================================================
%%  พรีแอมเบิลร่วม สำหรับสไลด์ประกอบการสอน บทที่ 1 : ตรรกศาสตร์ (Logic)
%%  รายวิชา 4091606 คณิตศาสตร์สำหรับคอมพิวเตอร์
%%  เรียงพิมพ์ด้วย XeLaTeX + ฟอนต์ TH Sarabun New
%%  ใช้ร่วมกันด้วย % !TeX program = xelatex
%% =====================================================================
%%  พรีแอมเบิลร่วม สำหรับสไลด์ประกอบการสอน บทที่ 1 : ตรรกศาสตร์ (Logic)
%%  รายวิชา 4091606 คณิตศาสตร์สำหรับคอมพิวเตอร์
%%  เรียงพิมพ์ด้วย XeLaTeX + ฟอนต์ TH Sarabun New
%%  ใช้ร่วมกันด้วย \input{_preamble} ที่บรรทัดแรกของไฟล์สไลด์แต่ละสัปดาห์
%% =====================================================================
%% โหมด handout: ไฟล์ไดรเวอร์ที่ \def\HANDOUTMODE{} ก่อน \input{_preamble}
%% จะรวมเลเยอร์ overlay (\pause) ให้แสดงเฉลยครบในหน้าเดียว (ไม่มีหน้าซ้ำ)
\ifdefined\HANDOUTMODE\PassOptionsToClass{handout}{beamer}\fi
\documentclass[aspectratio=169,11pt,xcolor=dvipsnames]{beamer}

\usepackage{fontspec}
\usepackage{amsmath,amssymb,amsthm}
\usepackage{booktabs}
\usepackage{array}
\usepackage{colortbl}
\usepackage{tikz}
\usetikzlibrary{arrows.meta,calc,positioning,shapes.gates.logic.US}

%% ---------- ภาษาไทย ----------
%% beamer ใช้แบบอักษร sans-serif เป็นค่าตั้งต้น จึงต้องตั้งทั้ง main และ sans
%% ให้เป็น TH Sarabun New และใช้ professionalfonts กันไม่ให้ beamer แทนที่ฟอนต์
\usefonttheme{professionalfonts}
\XeTeXlinebreaklocale "th_TH"
\XeTeXlinebreakskip = 0pt plus 1pt
\setmainfont[Script=Thai,Scale=1.25]{TH Sarabun New}
\setsansfont[Script=Thai,Scale=1.25]{TH Sarabun New}
\newfontfamily\thaifont[Script=Thai,Scale=1.25]{TH Sarabun New}

%% ---------- จานสี (ยึดตามตำรา) ----------
\definecolor{navy}{HTML}{1F3A93}
\definecolor{Tgreen}{HTML}{2E8B57}
\definecolor{Fred}{HTML}{C0392B}
\definecolor{gold}{HTML}{B7950B}
\definecolor{lightnavy}{HTML}{EAEEF7}
\definecolor{mint}{HTML}{E8F5E9}
\definecolor{lyellow}{HTML}{FFF6D5}
\definecolor{cream}{HTML}{FFF8E1}
\definecolor{lpink}{HTML}{FDEDEC}

%% ---------- ธีม / รูปลักษณ์ ----------
\usetheme{default}
\setbeamertemplate{navigation symbols}{}
\setbeamercolor{frametitle}{fg=white,bg=navy}
\setbeamercolor{title}{fg=navy}
\setbeamercolor{structure}{fg=navy}
\setbeamercolor{itemize item}{fg=navy}
\setbeamercolor{itemize subitem}{fg=gold}
\setbeamercolor{enumerate item}{fg=navy}
\setbeamercolor{section in toc}{fg=navy}
\setbeamercolor{block title}{bg=navy,fg=white}
\setbeamercolor{block body}{bg=lightnavy,fg=black}
\setbeamercolor{block title example}{bg=Tgreen,fg=white}
\setbeamercolor{block body example}{bg=mint,fg=black}
\setbeamercolor{block title alerted}{bg=Fred,fg=white}
\setbeamercolor{block body alerted}{bg=lpink,fg=black}
\setbeamerfont{frametitle}{series=\bfseries}
\setbeamerfont{title}{series=\bfseries}
\setbeamertemplate{itemize item}{\textbullet}
\setbeamertemplate{itemize subitem}{\textendash}

%% เส้นใต้หัวสไลด์
\setbeamertemplate{frametitle}{%
  \nointerlineskip
  \begin{beamercolorbox}[wd=\paperwidth,ht=2.6ex,dp=1.2ex,leftskip=0.3cm,rightskip=0.3cm]{frametitle}%
    \usebeamerfont{frametitle}\insertframetitle%
  \end{beamercolorbox}%
}

%% footline (กำหนดชื่อหัวเรื่องท้ายสไลด์ผ่าน \footchaptertitle)
\newcommand{\footchaptertitle}{บทที่ 1 ตรรกศาสตร์}
\setbeamertemplate{footline}{%
  \leavevmode%
  \hbox{%
    \begin{beamercolorbox}[wd=.5\paperwidth,ht=2.4ex,dp=1ex,leftskip=0.3cm]{author in head/foot}%
      \color{navy}\footnotesize คณิตศาสตร์สำหรับคอมพิวเตอร์ \,|\, \footchaptertitle%
    \end{beamercolorbox}%
    \begin{beamercolorbox}[wd=.5\paperwidth,ht=2.4ex,dp=1ex,rightskip=0.3cm]{author in head/foot}%
      \hfill\color{navy}\footnotesize \insertframenumber\,/\,\inserttotalframenumber%
    \end{beamercolorbox}%
  }%
}

%% ---------- มาโครช่วย ----------
\newcommand{\TT}{\textcolor{Tgreen}{\textbf{T}}}
\newcommand{\FF}{\textcolor{Fred}{\textbf{F}}}
\newcommand{\eng}[1]{\textcolor{gray}{\small (#1)}}
\renewcommand{\arraystretch}{1.25}

%% หน้าแบ่งหัวข้อ (section)
\AtBeginSection[]{%
  \begin{frame}[plain]
    \vfill
    \centering
    \begin{beamercolorbox}[sep=10pt,center,rounded=true,shadow=true]{title}
      \usebeamerfont{title}\Large\insertsectionhead
    \end{beamercolorbox}
    \vfill
  \end{frame}%
}
 ที่บรรทัดแรกของไฟล์สไลด์แต่ละสัปดาห์
%% =====================================================================
%% โหมด handout: ไฟล์ไดรเวอร์ที่ \def\HANDOUTMODE{} ก่อน % !TeX program = xelatex
%% =====================================================================
%%  พรีแอมเบิลร่วม สำหรับสไลด์ประกอบการสอน บทที่ 1 : ตรรกศาสตร์ (Logic)
%%  รายวิชา 4091606 คณิตศาสตร์สำหรับคอมพิวเตอร์
%%  เรียงพิมพ์ด้วย XeLaTeX + ฟอนต์ TH Sarabun New
%%  ใช้ร่วมกันด้วย \input{_preamble} ที่บรรทัดแรกของไฟล์สไลด์แต่ละสัปดาห์
%% =====================================================================
%% โหมด handout: ไฟล์ไดรเวอร์ที่ \def\HANDOUTMODE{} ก่อน \input{_preamble}
%% จะรวมเลเยอร์ overlay (\pause) ให้แสดงเฉลยครบในหน้าเดียว (ไม่มีหน้าซ้ำ)
\ifdefined\HANDOUTMODE\PassOptionsToClass{handout}{beamer}\fi
\documentclass[aspectratio=169,11pt,xcolor=dvipsnames]{beamer}

\usepackage{fontspec}
\usepackage{amsmath,amssymb,amsthm}
\usepackage{booktabs}
\usepackage{array}
\usepackage{colortbl}
\usepackage{tikz}
\usetikzlibrary{arrows.meta,calc,positioning,shapes.gates.logic.US}

%% ---------- ภาษาไทย ----------
%% beamer ใช้แบบอักษร sans-serif เป็นค่าตั้งต้น จึงต้องตั้งทั้ง main และ sans
%% ให้เป็น TH Sarabun New และใช้ professionalfonts กันไม่ให้ beamer แทนที่ฟอนต์
\usefonttheme{professionalfonts}
\XeTeXlinebreaklocale "th_TH"
\XeTeXlinebreakskip = 0pt plus 1pt
\setmainfont[Script=Thai,Scale=1.25]{TH Sarabun New}
\setsansfont[Script=Thai,Scale=1.25]{TH Sarabun New}
\newfontfamily\thaifont[Script=Thai,Scale=1.25]{TH Sarabun New}

%% ---------- จานสี (ยึดตามตำรา) ----------
\definecolor{navy}{HTML}{1F3A93}
\definecolor{Tgreen}{HTML}{2E8B57}
\definecolor{Fred}{HTML}{C0392B}
\definecolor{gold}{HTML}{B7950B}
\definecolor{lightnavy}{HTML}{EAEEF7}
\definecolor{mint}{HTML}{E8F5E9}
\definecolor{lyellow}{HTML}{FFF6D5}
\definecolor{cream}{HTML}{FFF8E1}
\definecolor{lpink}{HTML}{FDEDEC}

%% ---------- ธีม / รูปลักษณ์ ----------
\usetheme{default}
\setbeamertemplate{navigation symbols}{}
\setbeamercolor{frametitle}{fg=white,bg=navy}
\setbeamercolor{title}{fg=navy}
\setbeamercolor{structure}{fg=navy}
\setbeamercolor{itemize item}{fg=navy}
\setbeamercolor{itemize subitem}{fg=gold}
\setbeamercolor{enumerate item}{fg=navy}
\setbeamercolor{section in toc}{fg=navy}
\setbeamercolor{block title}{bg=navy,fg=white}
\setbeamercolor{block body}{bg=lightnavy,fg=black}
\setbeamercolor{block title example}{bg=Tgreen,fg=white}
\setbeamercolor{block body example}{bg=mint,fg=black}
\setbeamercolor{block title alerted}{bg=Fred,fg=white}
\setbeamercolor{block body alerted}{bg=lpink,fg=black}
\setbeamerfont{frametitle}{series=\bfseries}
\setbeamerfont{title}{series=\bfseries}
\setbeamertemplate{itemize item}{\textbullet}
\setbeamertemplate{itemize subitem}{\textendash}

%% เส้นใต้หัวสไลด์
\setbeamertemplate{frametitle}{%
  \nointerlineskip
  \begin{beamercolorbox}[wd=\paperwidth,ht=2.6ex,dp=1.2ex,leftskip=0.3cm,rightskip=0.3cm]{frametitle}%
    \usebeamerfont{frametitle}\insertframetitle%
  \end{beamercolorbox}%
}

%% footline (กำหนดชื่อหัวเรื่องท้ายสไลด์ผ่าน \footchaptertitle)
\newcommand{\footchaptertitle}{บทที่ 1 ตรรกศาสตร์}
\setbeamertemplate{footline}{%
  \leavevmode%
  \hbox{%
    \begin{beamercolorbox}[wd=.5\paperwidth,ht=2.4ex,dp=1ex,leftskip=0.3cm]{author in head/foot}%
      \color{navy}\footnotesize คณิตศาสตร์สำหรับคอมพิวเตอร์ \,|\, \footchaptertitle%
    \end{beamercolorbox}%
    \begin{beamercolorbox}[wd=.5\paperwidth,ht=2.4ex,dp=1ex,rightskip=0.3cm]{author in head/foot}%
      \hfill\color{navy}\footnotesize \insertframenumber\,/\,\inserttotalframenumber%
    \end{beamercolorbox}%
  }%
}

%% ---------- มาโครช่วย ----------
\newcommand{\TT}{\textcolor{Tgreen}{\textbf{T}}}
\newcommand{\FF}{\textcolor{Fred}{\textbf{F}}}
\newcommand{\eng}[1]{\textcolor{gray}{\small (#1)}}
\renewcommand{\arraystretch}{1.25}

%% หน้าแบ่งหัวข้อ (section)
\AtBeginSection[]{%
  \begin{frame}[plain]
    \vfill
    \centering
    \begin{beamercolorbox}[sep=10pt,center,rounded=true,shadow=true]{title}
      \usebeamerfont{title}\Large\insertsectionhead
    \end{beamercolorbox}
    \vfill
  \end{frame}%
}

%% จะรวมเลเยอร์ overlay (\pause) ให้แสดงเฉลยครบในหน้าเดียว (ไม่มีหน้าซ้ำ)
\ifdefined\HANDOUTMODE\PassOptionsToClass{handout}{beamer}\fi
\documentclass[aspectratio=169,11pt,xcolor=dvipsnames]{beamer}

\usepackage{fontspec}
\usepackage{amsmath,amssymb,amsthm}
\usepackage{booktabs}
\usepackage{array}
\usepackage{colortbl}
\usepackage{tikz}
\usetikzlibrary{arrows.meta,calc,positioning,shapes.gates.logic.US}

%% ---------- ภาษาไทย ----------
%% beamer ใช้แบบอักษร sans-serif เป็นค่าตั้งต้น จึงต้องตั้งทั้ง main และ sans
%% ให้เป็น TH Sarabun New และใช้ professionalfonts กันไม่ให้ beamer แทนที่ฟอนต์
\usefonttheme{professionalfonts}
\XeTeXlinebreaklocale "th_TH"
\XeTeXlinebreakskip = 0pt plus 1pt
\setmainfont[Script=Thai,Scale=1.25]{TH Sarabun New}
\setsansfont[Script=Thai,Scale=1.25]{TH Sarabun New}
\newfontfamily\thaifont[Script=Thai,Scale=1.25]{TH Sarabun New}

%% ---------- จานสี (ยึดตามตำรา) ----------
\definecolor{navy}{HTML}{1F3A93}
\definecolor{Tgreen}{HTML}{2E8B57}
\definecolor{Fred}{HTML}{C0392B}
\definecolor{gold}{HTML}{B7950B}
\definecolor{lightnavy}{HTML}{EAEEF7}
\definecolor{mint}{HTML}{E8F5E9}
\definecolor{lyellow}{HTML}{FFF6D5}
\definecolor{cream}{HTML}{FFF8E1}
\definecolor{lpink}{HTML}{FDEDEC}

%% ---------- ธีม / รูปลักษณ์ ----------
\usetheme{default}
\setbeamertemplate{navigation symbols}{}
\setbeamercolor{frametitle}{fg=white,bg=navy}
\setbeamercolor{title}{fg=navy}
\setbeamercolor{structure}{fg=navy}
\setbeamercolor{itemize item}{fg=navy}
\setbeamercolor{itemize subitem}{fg=gold}
\setbeamercolor{enumerate item}{fg=navy}
\setbeamercolor{section in toc}{fg=navy}
\setbeamercolor{block title}{bg=navy,fg=white}
\setbeamercolor{block body}{bg=lightnavy,fg=black}
\setbeamercolor{block title example}{bg=Tgreen,fg=white}
\setbeamercolor{block body example}{bg=mint,fg=black}
\setbeamercolor{block title alerted}{bg=Fred,fg=white}
\setbeamercolor{block body alerted}{bg=lpink,fg=black}
\setbeamerfont{frametitle}{series=\bfseries}
\setbeamerfont{title}{series=\bfseries}
\setbeamertemplate{itemize item}{\textbullet}
\setbeamertemplate{itemize subitem}{\textendash}

%% เส้นใต้หัวสไลด์
\setbeamertemplate{frametitle}{%
  \nointerlineskip
  \begin{beamercolorbox}[wd=\paperwidth,ht=2.6ex,dp=1.2ex,leftskip=0.3cm,rightskip=0.3cm]{frametitle}%
    \usebeamerfont{frametitle}\insertframetitle%
  \end{beamercolorbox}%
}

%% footline (กำหนดชื่อหัวเรื่องท้ายสไลด์ผ่าน \footchaptertitle)
\newcommand{\footchaptertitle}{บทที่ 1 ตรรกศาสตร์}
\setbeamertemplate{footline}{%
  \leavevmode%
  \hbox{%
    \begin{beamercolorbox}[wd=.5\paperwidth,ht=2.4ex,dp=1ex,leftskip=0.3cm]{author in head/foot}%
      \color{navy}\footnotesize คณิตศาสตร์สำหรับคอมพิวเตอร์ \,|\, \footchaptertitle%
    \end{beamercolorbox}%
    \begin{beamercolorbox}[wd=.5\paperwidth,ht=2.4ex,dp=1ex,rightskip=0.3cm]{author in head/foot}%
      \hfill\color{navy}\footnotesize \insertframenumber\,/\,\inserttotalframenumber%
    \end{beamercolorbox}%
  }%
}

%% ---------- มาโครช่วย ----------
\newcommand{\TT}{\textcolor{Tgreen}{\textbf{T}}}
\newcommand{\FF}{\textcolor{Fred}{\textbf{F}}}
\newcommand{\eng}[1]{\textcolor{gray}{\small (#1)}}
\renewcommand{\arraystretch}{1.25}

%% หน้าแบ่งหัวข้อ (section)
\AtBeginSection[]{%
  \begin{frame}[plain]
    \vfill
    \centering
    \begin{beamercolorbox}[sep=10pt,center,rounded=true,shadow=true]{title}
      \usebeamerfont{title}\Large\insertsectionhead
    \end{beamercolorbox}
    \vfill
  \end{frame}%
}

%% จะรวมเลเยอร์ overlay (\pause) ให้แสดงเฉลยครบในหน้าเดียว (ไม่มีหน้าซ้ำ)
\ifdefined\HANDOUTMODE\PassOptionsToClass{handout}{beamer}\fi
\documentclass[aspectratio=169,11pt,xcolor=dvipsnames]{beamer}

\usepackage{fontspec}
\usepackage{amsmath,amssymb,amsthm}
\usepackage{booktabs}
\usepackage{array}
\usepackage{colortbl}
\usepackage{tikz}
\usetikzlibrary{arrows.meta,calc,positioning,shapes.gates.logic.US}

%% ---------- ภาษาไทย ----------
%% beamer ใช้แบบอักษร sans-serif เป็นค่าตั้งต้น จึงต้องตั้งทั้ง main และ sans
%% ให้เป็น TH Sarabun New และใช้ professionalfonts กันไม่ให้ beamer แทนที่ฟอนต์
\usefonttheme{professionalfonts}
\XeTeXlinebreaklocale "th_TH"
\XeTeXlinebreakskip = 0pt plus 1pt
\setmainfont[Script=Thai,Scale=1.25]{TH Sarabun New}
\setsansfont[Script=Thai,Scale=1.25]{TH Sarabun New}
\newfontfamily\thaifont[Script=Thai,Scale=1.25]{TH Sarabun New}

%% ---------- จานสี (ยึดตามตำรา) ----------
\definecolor{navy}{HTML}{1F3A93}
\definecolor{Tgreen}{HTML}{2E8B57}
\definecolor{Fred}{HTML}{C0392B}
\definecolor{gold}{HTML}{B7950B}
\definecolor{lightnavy}{HTML}{EAEEF7}
\definecolor{mint}{HTML}{E8F5E9}
\definecolor{lyellow}{HTML}{FFF6D5}
\definecolor{cream}{HTML}{FFF8E1}
\definecolor{lpink}{HTML}{FDEDEC}

%% ---------- ธีม / รูปลักษณ์ ----------
\usetheme{default}
\setbeamertemplate{navigation symbols}{}
\setbeamercolor{frametitle}{fg=white,bg=navy}
\setbeamercolor{title}{fg=navy}
\setbeamercolor{structure}{fg=navy}
\setbeamercolor{itemize item}{fg=navy}
\setbeamercolor{itemize subitem}{fg=gold}
\setbeamercolor{enumerate item}{fg=navy}
\setbeamercolor{section in toc}{fg=navy}
\setbeamercolor{block title}{bg=navy,fg=white}
\setbeamercolor{block body}{bg=lightnavy,fg=black}
\setbeamercolor{block title example}{bg=Tgreen,fg=white}
\setbeamercolor{block body example}{bg=mint,fg=black}
\setbeamercolor{block title alerted}{bg=Fred,fg=white}
\setbeamercolor{block body alerted}{bg=lpink,fg=black}
\setbeamerfont{frametitle}{series=\bfseries}
\setbeamerfont{title}{series=\bfseries}
\setbeamertemplate{itemize item}{\textbullet}
\setbeamertemplate{itemize subitem}{\textendash}

%% เส้นใต้หัวสไลด์
\setbeamertemplate{frametitle}{%
  \nointerlineskip
  \begin{beamercolorbox}[wd=\paperwidth,ht=2.6ex,dp=1.2ex,leftskip=0.3cm,rightskip=0.3cm]{frametitle}%
    \usebeamerfont{frametitle}\insertframetitle%
  \end{beamercolorbox}%
}

%% footline (กำหนดชื่อหัวเรื่องท้ายสไลด์ผ่าน \footchaptertitle)
\newcommand{\footchaptertitle}{บทที่ 1 ตรรกศาสตร์}
\setbeamertemplate{footline}{%
  \leavevmode%
  \hbox{%
    \begin{beamercolorbox}[wd=.5\paperwidth,ht=2.4ex,dp=1ex,leftskip=0.3cm]{author in head/foot}%
      \color{navy}\footnotesize คณิตศาสตร์สำหรับคอมพิวเตอร์ \,|\, \footchaptertitle%
    \end{beamercolorbox}%
    \begin{beamercolorbox}[wd=.5\paperwidth,ht=2.4ex,dp=1ex,rightskip=0.3cm]{author in head/foot}%
      \hfill\color{navy}\footnotesize \insertframenumber\,/\,\inserttotalframenumber%
    \end{beamercolorbox}%
  }%
}

%% ---------- มาโครช่วย ----------
\newcommand{\TT}{\textcolor{Tgreen}{\textbf{T}}}
\newcommand{\FF}{\textcolor{Fred}{\textbf{F}}}
\newcommand{\eng}[1]{\textcolor{gray}{\small (#1)}}
\renewcommand{\arraystretch}{1.25}

%% หน้าแบ่งหัวข้อ (section)
\AtBeginSection[]{%
  \begin{frame}[plain]
    \vfill
    \centering
    \begin{beamercolorbox}[sep=10pt,center,rounded=true,shadow=true]{title}
      \usebeamerfont{title}\Large\insertsectionhead
    \end{beamercolorbox}
    \vfill
  \end{frame}%
}

%% จะรวมเลเยอร์ overlay (\pause) ให้แสดงเฉลยครบในหน้าเดียว (ไม่มีหน้าซ้ำ)
\ifdefined\HANDOUTMODE\PassOptionsToClass{handout}{beamer}\fi
\documentclass[aspectratio=169,11pt,xcolor=dvipsnames]{beamer}

\usepackage{fontspec}
\usepackage{amsmath,amssymb,amsthm}
\usepackage{booktabs}
\usepackage{array}
\usepackage{colortbl}
\usepackage{tikz}
\usetikzlibrary{arrows.meta,calc,positioning,shapes.gates.logic.US}

%% ---------- ภาษาไทย ----------
%% beamer ใช้แบบอักษร sans-serif เป็นค่าตั้งต้น จึงต้องตั้งทั้ง main และ sans
%% ให้เป็น TH Sarabun New และใช้ professionalfonts กันไม่ให้ beamer แทนที่ฟอนต์
\usefonttheme{professionalfonts}
\XeTeXlinebreaklocale "th_TH"
\XeTeXlinebreakskip = 0pt plus 1pt
\setmainfont[Script=Thai,Scale=1.25]{TH Sarabun New}
\setsansfont[Script=Thai,Scale=1.25]{TH Sarabun New}
\newfontfamily\thaifont[Script=Thai,Scale=1.25]{TH Sarabun New}

%% ---------- จานสี (ยึดตามตำรา) ----------
\definecolor{navy}{HTML}{1F3A93}
\definecolor{Tgreen}{HTML}{2E8B57}
\definecolor{Fred}{HTML}{C0392B}
\definecolor{gold}{HTML}{B7950B}
\definecolor{lightnavy}{HTML}{EAEEF7}
\definecolor{mint}{HTML}{E8F5E9}
\definecolor{lyellow}{HTML}{FFF6D5}
\definecolor{cream}{HTML}{FFF8E1}
\definecolor{lpink}{HTML}{FDEDEC}

%% ---------- ธีม / รูปลักษณ์ ----------
\usetheme{default}
\setbeamertemplate{navigation symbols}{}
\setbeamercolor{frametitle}{fg=white,bg=navy}
\setbeamercolor{title}{fg=navy}
\setbeamercolor{structure}{fg=navy}
\setbeamercolor{itemize item}{fg=navy}
\setbeamercolor{itemize subitem}{fg=gold}
\setbeamercolor{enumerate item}{fg=navy}
\setbeamercolor{section in toc}{fg=navy}
\setbeamercolor{block title}{bg=navy,fg=white}
\setbeamercolor{block body}{bg=lightnavy,fg=black}
\setbeamercolor{block title example}{bg=Tgreen,fg=white}
\setbeamercolor{block body example}{bg=mint,fg=black}
\setbeamercolor{block title alerted}{bg=Fred,fg=white}
\setbeamercolor{block body alerted}{bg=lpink,fg=black}
\setbeamerfont{frametitle}{series=\bfseries}
\setbeamerfont{title}{series=\bfseries}
\setbeamertemplate{itemize item}{\textbullet}
\setbeamertemplate{itemize subitem}{\textendash}

%% เส้นใต้หัวสไลด์
\setbeamertemplate{frametitle}{%
  \nointerlineskip
  \begin{beamercolorbox}[wd=\paperwidth,ht=2.6ex,dp=1.2ex,leftskip=0.3cm,rightskip=0.3cm]{frametitle}%
    \usebeamerfont{frametitle}\insertframetitle%
  \end{beamercolorbox}%
}

%% footline (กำหนดชื่อหัวเรื่องท้ายสไลด์ผ่าน \footchaptertitle)
\newcommand{\footchaptertitle}{บทที่ 1 ตรรกศาสตร์}
\setbeamertemplate{footline}{%
  \leavevmode%
  \hbox{%
    \begin{beamercolorbox}[wd=.5\paperwidth,ht=2.4ex,dp=1ex,leftskip=0.3cm]{author in head/foot}%
      \color{navy}\footnotesize คณิตศาสตร์สำหรับคอมพิวเตอร์ \,|\, \footchaptertitle%
    \end{beamercolorbox}%
    \begin{beamercolorbox}[wd=.5\paperwidth,ht=2.4ex,dp=1ex,rightskip=0.3cm]{author in head/foot}%
      \hfill\color{navy}\footnotesize \insertframenumber\,/\,\inserttotalframenumber%
    \end{beamercolorbox}%
  }%
}

%% ---------- มาโครช่วย ----------
\newcommand{\TT}{\textcolor{Tgreen}{\textbf{T}}}
\newcommand{\FF}{\textcolor{Fred}{\textbf{F}}}
\newcommand{\eng}[1]{\textcolor{gray}{\small (#1)}}
\renewcommand{\arraystretch}{1.25}

%% หน้าแบ่งหัวข้อ (section)
\AtBeginSection[]{%
  \begin{frame}[plain]
    \vfill
    \centering
    \begin{beamercolorbox}[sep=10pt,center,rounded=true,shadow=true]{title}
      \usebeamerfont{title}\Large\insertsectionhead
    \end{beamercolorbox}
    \vfill
  \end{frame}%
}
